\documentclass[12pt]{article}
\usepackage[utf8]{inputenc}
\usepackage[T1]{fontenc}
\usepackage{hyperref}
\title{Patep Language Resources for K-12 Curriculum Development}
\author{Compiled by Claude Sonnet 5.5 (claude-sonnet-5-5) for Andr\'as Kornai}
\date{}
\begin{document}
\maketitle

\section{Introduction}
Patep (ISO 639-3: ptp; Glottocode pate1247; also Ptep, Dengalu) is an Austronesian language of Morobe Province, Papua New Guinea, in the Mumeng dialect chain (with Zenag and Gorakor). The search results give about 1,700 speakers in the Mumeng area, far fewer than the roughly 10,000 in the task description; this discrepancy should be checked. Little material specific to Patep was found.

\section{Primary Educational Resources}
No primers, textbooks or school programmes in Patep were found in the searches.

\section{Linguistic and Reference Materials}
\begin{itemize}
\item Wikipedia article on the Mumeng language (dialect chain including Patep): \url{https://en.wikipedia.org/wiki/Mumeng_language}
\item Counting-system record for Patep (Max Planck EVA, Numeral Systems of Papua New Guinea): \url{https://lingweb.eva.mpg.de/channumerals/Patep.htm}
\item Austronesian Comparative Dictionary language page: \url{https://acd.clld.org/languages/923}
\end{itemize}

\section{Religious and Community Resources}
\begin{itemize}
\item Joshua Project profile of Patep in Papua New Guinea, which reports that a New Testament has been translated: \url{https://joshuaproject.net/people_groups/14337}
\end{itemize}

\section{Digital and Audio Resources}
Global Recordings Network has a Patep language page, but its listing indicates no recordings are available: \url{https://globalrecordings.net/en/language/ptp}

\section{Curriculum Development Context}
Patep is not known to be taught in schools. The searches found only overview and reference material, and a reported New Testament translation which could be a source of written text if obtained from the translating agency. Curriculum development would require new materials produced with the community.

\end{document}